\FloatBarrier
\section{Related Work}
\label{sec:related_work}
Our work targets \emph{context-grounded, tool-using} agents for time-series analysis.
We situate \name among \textit{(i)} data-analysis and ML-agent benchmarks, \textit{(ii)} time-series reasoning and temporal-textual question-answer benchmarks, and \textit{(iii)} agent benchmarks that primarily measure end-to-end code execution.
\Cref{tab:benchmark_comparison} summarizes how closely related time-series benchmarks differ along dimensions most relevant to our goals, including controllability of \emph{text} and \emph{numeric} context and whether they evaluate \emph{agentic execution}.

\begin{table*}[t]
    \centering
    \footnotesize
    \setlength{\tabcolsep}{6pt}
    \renewcommand{\arraystretch}{1.12}
    \caption{Comparison with the most relevant benchmarks.
    ``Controlled text context'' indicates whether the benchmark provides multiple versions of the same task with systematically varied levels of natural-language context.
    ``Controlled numerical context'' indicates whether the benchmark provides multiple versions of the same task with systematically varied exemplars or levels of access to the time series.
    ``uni'' stands for univariate, ``multi'' for multivariate.}
    \label{tab:benchmark_comparison}
    \NoHyper
    \resizebox{\linewidth}{!}{
    \begin{tabular}{lcccccc}
    \toprule
    \textbf{Benchmark} &
    \makecell{\textbf{Synthetic}\\\textbf{xata}} &
    \makecell{\textbf{Controlled}\\\textbf{text context}} &
    \makecell{\textbf{Controlled}\\\textbf{numeric context}} &
    \makecell{\textbf{Time Series}\\\textbf{type}} &
    \makecell{\textbf{Agentic}\\\textbf{execution}} &
    \makecell{\textbf{Context in}\\\textbf{prompt}} \\
    \midrule
    \makecell[l]{TimeSeriesExamAgent\\\citep{gwiazda2026timeseriesexamagent}}
    & \makecell{LLM-generated}
    & $\times$
    & $\times$
    & uni
    & $\times$
    & $\checkmark$ \\
    \makecell[l]{TSRBench\\\citep{yu2026tsrbench}}
    & \makecell{Domain-informed\\Code-generation}
    & $\times$
    & $\times$
    & multi
    & $\times$
    & $\checkmark$ \\
    \makecell[l]{MTBench\\\citep{chen2025mtbench}}
    & \makecell{LLM-assisted\\Annotation/Alignment}
    & $\times$
    & $\times$
    & uni
    & $\times$
    & $\checkmark$ \\
    \makecell[l]{Language Models\\\citep{merrill2024language}}
    & \makecell{LLM-generated}
    & $\times$
    & $\times$
    & uni
    & $\times$
    & $\checkmark$ \\
    \midrule
    \makecell[l]{\textbf{\name\ (ours)}}
    & \makecell{Simulation-based,\\Physical ground truth}
    & $\checkmark$
    & $\checkmark$
    & multi
    & \makecell{$\checkmark$}
    & \makecell{$\checkmark$} \\
    \bottomrule
    \end{tabular}
    }
    \endNoHyper
\end{table*}

\subsection{Hybrid text--data reasoning and data-analysis benchmarks}
Early hybrid QA benchmarks such as FinQA~\citep{chen2021finqa} and TAT-QA~\citep{zhu2021tatqa} combine natural-language passages with structured tables and require multi-step numerical reasoning.
These datasets helped establish evaluation protocols for \emph{text + numbers} reasoning, but they are typically single-step QA and do not require an agent to decide \emph{what evidence to gather} from raw data.

More recent work evaluates data-analysis workflows with code execution.
InfiAgent-DABench~\citep{hu2024infidabench} benchmarks tool-using agents in an execution environment, while InsightBench~\citep{sahu2024insightbench} evaluates multi-step business analytics scenarios.
DSBench~\citep{jing2024dsbench} and DACO~\citep{wu2024daco} further stress end-to-end data-science pipelines and code generation, and systems such as DS-Agent~\citep{guo2024ds} explore case-based reasoning for automated data science.
While these benchmarks probe workflow completion, the relevant ``context'' is often sparse or directly encoded in schemas and dataset descriptions, and tasks can frequently be solved by applying generic procedures without validating assumptions against raw values.

\subsection{Time-series reasoning, forecasting, and temporal-textual QA}
Several benchmarks probe whether LLMs understand core time-series concepts.
TimeSeriesExam~\citep{cai2024timeseriesexam} provides a scalable multiple-choice exam across pattern recognition, noise, similarity, anomaly detection, and causality, but does not evaluate interactive exploration.
\citet{merrill2024language} show that strong language models still struggle to reason about time series when signals are presented as raw numeric sequences, motivating tool-using and multimodal approaches.

In parallel, a line of work studies time-series forecasting and reasoning with LLMs.
Gruver et al.~\citep{gruver2024large} demonstrate that LLMs can act as zero-shot time-series forecasters, while Tan et al.~\citep{tan2024are} and Chow et al.~\citep{chow2024towards} analyze when language-model priors help or hurt and propose interfaces for time-series reasoning.
SciTS~\citep{anonymous2025scits} studies scientific time-series understanding and generation with LLMs.
These settings typically expose the time series through a fixed representation and evaluate single-pass answers, rather than testing whether an \emph{agent} can iteratively choose which segment, statistic, or diagnostic to inspect to ground its decision.

Temporal-textual QA resources couple time-series signals with language.
MTBench~\citep{chen2025mtbench} evaluates models on joint reasoning over paired time-series and text, and ITFormer~\citep{wang2025itformer} proposes datasets and model interfaces for temporal-textual QA.
ChatTS~\citep{xie2024chatts} demonstrates that synthetic alignment data can substantially improve multimodal understanding and reasoning about time series.
However, these benchmarks do not isolate whether decisions are grounded in evidence gathered through targeted, multi-step exploration of raw values.

\subsection{From code-centric agents to evidence-grounded time-series agents}
Agent benchmarks such as MLE-bench~\citep{chan2024mle} evaluate end-to-end ML engineering skills in realistic workflows, emphasizing the ability to write and execute code.
Recent evidence suggests that time-series domains remain challenging even for capable agents: Cai et al.~\citep{cai2025llmAgents} find that agents often fail to produce reliable time-series solutions despite competent tool use.
Such benchmarks are valuable for measuring pipeline completion, but aggregate outcome metrics can obscure if the agent's decisions are supported by the observed temporal dynamics.

\name complements this line of work by targeting a narrower, reasoning-centric regime: tasks are constructed so that \textit{(i)} natural-language context is necessary but not sufficient, and \textit{(ii)} correct solutions require targeted inspection of the raw signal under potentially misleading priors.
By combining simulation episodes with known ground-truth interventions and controlled context ablations, \name is designed to diagnose whether agents \emph{ground} their conclusions in time-series evidence rather than relying on default pipelines or textual shortcuts.
